
\begin{figure}
%\centerline{\includegraphics[width=9cm,height=6cm]{images/scope1.eps}}
\setlength{\unitlength}{1cm}
\begin{picture}(10,3)
\linethickness{0.3mm}

\put(0.1,2.6){\large{Command line: \codefont{>a.out these are args}}}
\put(0.1,2.0){\large{Resulting data structures:}}
\put(4.3,1.5){\codefont{argv[0]} \hspace*{1.5cm}  a . o u t \textbackslash 0}
\put(4.3,1.0){\codefont{argv[1]}  \hspace*{1.5cm} t h e s e \textbackslash 0}
\put(4.3,0.5){\codefont{argv[2]}  \hspace*{1.5cm} a r e \textbackslash 0}
\put(4.3,0.0){\codefont{argv[3]}  \hspace*{1.5cm} a r g s \textbackslash 0}
\put(5.8,1.6){\vector(1,0){1.3}}
\put(5.8,1.1){\vector(1,0){1.3}}
\put(5.8,0.6){\vector(1,0){1.3}}
\put(5.8,0.1){\vector(1,0){1.3}}

\put(1.5,1.5){\codefont{argc}}
\put(2.6,1.0){\codefont{4}}
\put(2.5,0.9){\line(0,1){0.5}}
\put(3.5,0.9){\line(0,1){0.5}}
\put(2.5,0.9){\line(1,0){1.0}}
\put(2.5,1.4){\line(1,0){1.0}}

\end{picture}
\caption{Illustration of command line arguments in C++.  This shows the data structures created when the command line command: ``\codefont{a.out these are args}" is run.  The program name is \codefont{a.out}, the command line arguments are:  ``these'', ``are'', and ``args''. The parameter \codefont{argc} gets the value 4 because there are 4 arguments total, including the program name (\codefont{a.out} in this example), which is also treated as an argument.  Note that \codefont{argv} is an array of pointers to C-style strings and each argument is stored as a C-style array of characters (an array of individual characters ending with the \codefont{null} character \textbackslash 0.}
\label{fig:args}
\end{figure}

Because the \codefont{main()} function is a function it can both return a value and can accept arguments.  The projects in this text don't involve passing arguments to \codefont{main()} because such arguments are generally included by running the program, with arguments, from a command line and many programming environments and IDEs by default don't run programs from the command line.  It is always possible to run a program from the command line, but the steps vary somewhat by IDE and operating system.

As an example, in a Windows environment running a program from the command requires opening  a command window, moving to the directory containing the compiled program (the executable), and running the program with associated arguments.  Running a program called \codefont{a.out} from the command line with arguments would look like:\\
\codefont{
>a.out these are args\\
}
where \codefont{>} is the command line prompt and the extra arguments are ``these'', ``are'', and ``args''.  By default the name of the program, in this case ``a.out'' is also passed to the program as an argument.  This allows a program to `know' its own name when it runs.

To write a C++ program that accepts command line arguments the \codefont{main()} function is written as:\\
\codefont{
int main(int argc, char* argv[])\footnote{Another, equivalent option is \codefont{int main(int argc, char **argv)}}\\
}
with exactly those variable names.  The variable \codefont{argc} (short for ARGument Count) is the number of arguments passed to \codefont{main()}.  The arguments passed to a C++ program  include the name of the program itself.  The variable \codefont{argv} (short for ARGument Vector) is an array of pointers to C-style character strings.   Figure~\ref{fig:args} illustrates these data structures and their contents if the program is run using the arguments listed above.  Notice that the first string pointed to by the \codefont{argv} array is the program name \codefont{a.out}.

It is important to notice that the arguments to \codefont{main()} are \emph{always} treated as strings, even if they are meant to be integers.  E.g. if a program is run with the command line argument \codefont{415} it is treated as a string of characters: `4', `1', `5'.  To get the argument as an actual integer a function such as \codefont{atoi()} needs to be used.  The function \codefont{atoi()} is defined in the \codefont{cstdlib}, is short for Ascii TO Integer and converts a C-style string into an integer.  E.g. the command:\\
\codefont{X = atoi(``415'');}\\
will change the string ``415'' into the integer \codefont{415} and store it in \codefont{X}.

Command line arguments can be very useful for setting important parameters within a program.  For example, in the NIM program a command line argument could be used to set the initial number of objects.  In the Pet program a command line argument could be used to supply the name of the pet to be played with (or the name of the file containing the pet if you've modified the program to save and reload pets).  In the Generic Board Game a command line argument could be used to supply the name of the game to be played (again really this would be the name of the game file to open).  Of course, all of this data could also be supplied using standard I/O commands when the program is running.  Command line arguments are generally only \emph{required} when one program is starting another program and can't easily respond to user oriented I/O.